\section{Realios mašinos projektas}

\subsection{Techninės įrangos komponentai}

\begin{description}
\item[Procesorius] gali dirbti dviem režimais – supervizoriaus arba vartotojo. Naudotojo rėžime vykdomos virtualios mašinos komandos, o supervizoriaus rėžime – operacinės sistemos. Laiko tarpas, kurį procesorius atlieka vieną komandą laikomas vienu laiko vienetu. Procesorius turi šiuos registrus:
    \begin{description}
        \item[$IC$] 2 baitų nuorodos registras. Skirtas nurodyti vykdomos 
        virtualios mašinos komandos adresą atmintyje.
        \item[$R1$] 4 baitų bendro naudojimo registras.
        \item[$R2$] 4 baitų bendro naudojimo registras.
        \item[$SF$] 2 baitų loginis registras, saugomi aritmetinių ir loginių operacijų suformuoti požymiai.
        \item[$MODE$] 1 baito loginis registras, kurio reikšmė nusako 
        procesoriaus darbo rėžimą („N“ – naudotojas, „S“ – supervizorius).
    \end{description}
\end{description}

        